% BHU PYQ -- Transcribed & Corrected
% Paper: ENGAE-11 / ENGAE11: Writing Skills
% Exam: Under Graduate Semester I Examination, 2024-25 / 2025-26 (NEP)
% Category: Ability Enhancement Course (AEC)

\documentclass[11pt,a4paper]{article}
\usepackage[a4paper,top=2.5cm,bottom=2.5cm,left=2.8cm,right=2.8cm,headheight=14pt]{geometry}
\usepackage{amsmath,amssymb}
\usepackage[T1]{fontenc}
\usepackage[utf8]{inputenc}
\usepackage{lmodern,microtype,fancyhdr,enumitem,hyperref,lastpage,parskip}

\hypersetup{
    colorlinks=true,
    urlcolor=black,
    linkcolor=black,
    pdftitle={ENGAE11: Writing Skills -- BHU Sem I 2024-25 (NEP AEC)}
}

\pagestyle{fancy}
\fancyhf{}
\renewcommand{\headrulewidth}{0.4pt}
\renewcommand{\footrulewidth}{0.4pt}
\fancyhead[L]{\small   \textit{Banaras Hindu University}}
\fancyhead[R]{\small   \textit{ENGAE-11: Writing Skills (AEC)}}
\fancyfoot[C]{\small Page  \thepage\ of \pageref{LastPage}}


\newlist{parts}{enumerate}{2}
\setlist[parts,1]{label=(\Alph*),leftmargin=*,itemsep=3pt,topsep=2pt}

\newcommand{\pts}[1]{\hfill   \textit{[#1]}}


\begin{document}


\begin{center}
  {\large\bfseries BANARAS HINDU UNIVERSITY}\\[2pt]
  {\normalsize Under Graduate --- Semester I Examination, 2024-25 / 2025-26 (NEP)}\\[6pt]
  \rule{0.8\linewidth}{0.5pt}\\[6pt]
  {\large\bfseries ABILITY ENHANCEMENT COURSE (AEC)}\\[4pt]
  {\large\bfseries Paper Code: ENGAE-11 : Writing Skills}\\[4pt]
  {\normalsize Time: 1.5 Hours\hfill Maximum Marks: 70}
\end{center}

\rule{\linewidth}{0.4pt}

\smallskip



\noindent   \textit{Instructions: Answer all 70 Multiple Choice Questions. Each question carries equal marks (1 mark each).}

\bigskip




\noindent   \textbf{Question 1.} What should the conclusion of an essay primarily do?\pts{1}

\begin{parts}
  \item Summarize the main points and restate the thesis in a new way
  \item Introduce new arguments
  \item Ask unrelated questions
  \item End abruptly without summary
\end{parts}

\medskip




\noindent   \textbf{Question 2.} Which of the following is the correct sequence of story development?\pts{1}

\begin{parts}
  \item Conflict $ o$ Ending $ o$ Climax $ o$ Beginning
  \item Beginning $ o$ Conflict $ o$ Climax $ o$ Ending
  \item Climax $ o$ Beginning $ o$ Conflict $ o$ Ending
  \item Ending $ o$ Conflict $ o$ Climax $ o$ Beginning
\end{parts}

\medskip




\noindent   \textbf{Question 3.} In job application letters, which paragraph typically outlines your skills and experiences?\pts{1}

\begin{parts}
  \item Salutation
  \item Body paragraph(s)
  \item Closing
  \item Subject line
\end{parts}

\medskip




\noindent   \textbf{Question 4.} What is the main purpose of a cover letter?\pts{1}

\begin{parts}
  \item To introduce yourself and explain why you are qualified for the job
  \item To list your entire life history
  \item To summarize your financial status
  \item To ask for personal favors
\end{parts}

\medskip




\noindent   \textbf{Question 5.} What is the purpose of the ``CC'' field in an email?\pts{1}

\begin{parts}
  \item To send an email to a friend
  \item To add recipients who should be informed but do not need to take action
  \item To send an email to your neighbor
  \item To mark an email as important
\end{parts}

\medskip




\noindent   \textbf{Question 6.} The process of revising an essay typically includes which of the following?\pts{1}

\begin{parts}
  \item Checking for grammar, clarity, and logical coherence
  \item Adding unrelated topics
  \item Reducing word count without purpose
  \item Changing main ideas repeatedly
\end{parts}

\medskip




\noindent   \textbf{Question 7.} Which of the following is an example of an official/formal letter?\pts{1}

\begin{parts}
  \item Letter to a business manager or university authority
  \item Letter to a sibling
  \item Birthday invitation to a friend
  \item Personal diary entry
\end{parts}

\medskip




\noindent   \textbf{Question 8.} In email writing, what does ``BCC'' stand for?\pts{1}

\begin{parts}
  \item Blind Carbon Copy
  \item Basic Carbon Copy
  \item Blind Courtesy Communication
  \item Business Communication Copy
\end{parts}

\medskip




\noindent   \textbf{Question 9.} What is the primary focus of Précis writing?\pts{1}

\begin{parts}
  \item Condensing a text to its essential ideas without altering the original meaning
  \item Adding personal opinions and commentary
  \item Expanding the length of the original text
  \item Changing the author's viewpoint
\end{parts}

\medskip




\noindent   \textbf{Question 10.} What is the ideal length of a Précis relative to the original passage?\pts{1}

\begin{parts}
  \item Approximately one-third of the original passage
  \item Half of the original passage
  \item Same length as the original passage
  \item Double the length of the original passage
\end{parts}

\medskip




\noindent   \textbf{Question 11.} Which tone is appropriate for formal business letters?\pts{1}

\begin{parts}
  \item Polite, professional, and clear
  \item Sarcastic and informal
  \item Overly emotional and dramatic
  \item Casual and humorous
\end{parts}

\medskip




\noindent   \textbf{Question 12.} What element should be included in the header of a formal letter?\pts{1}

\begin{parts}
  \item Sender's address, date, and receiver's address
  \item Informal greeting only
  \item Personal hobbies
  \item Social media handles
\end{parts}

\medskip




\noindent   \textbf{Question 13.} In essay writing, what is a thesis statement?\pts{1}

\begin{parts}
  \item A single sentence that summarizes the main point or claim of the essay
  \item A list of references
  \item The concluding sentence of a paragraph
  \item An introductory greeting
\end{parts}

\medskip




\noindent   \textbf{Question 14.} What type of essay aims to convince the reader to accept a specific viewpoint?\pts{1}

\begin{parts}
  \item Persuasive / Argumentative Essay
  \item Descriptive Essay
  \item Narrative Essay
  \item Expository Essay
\end{parts}

\medskip




\noindent   \textbf{Question 15.} What type of essay describes a person, place, object, or event vividly?\pts{1}

\begin{parts}
  \item Descriptive Essay
  \item Persuasive Essay
  \item Argumentative Essay
  \item Analytical Essay
\end{parts}

\medskip




\noindent   \textbf{Question 16.} What is the first step in the writing process?\pts{1}

\begin{parts}
  \item Pre-writing / Brainstorming
  \item Proofreading
  \item Publishing
  \item Final drafting
\end{parts}

\medskip




\noindent   \textbf{Question 17.} Which salutation is most appropriate for a formal letter when the recipient's name is unknown?\pts{1}

\begin{parts}
  \item Dear Sir / Madam
  \item Hey friend
  \item To Whom It May Concern
  \item Both (A) and (C)
\end{parts}

\medskip




\noindent   \textbf{Question 18.} What is considered inappropriate in professional email communication?\pts{1}

\begin{parts}
  \item Using ALL CAPS (which implies shouting) and informal slang
  \item Clear subject line
  \item Professional closing
  \item Checking spelling before sending
\end{parts}

\medskip




\noindent   \textbf{Question 19.} What should be avoided when writing a précis?\pts{1}

\begin{parts}
  \item Inserting personal opinions, extra examples, or direct quotes
  \item Maintaining original meaning
  \item Concise summary
  \item Using clear language
\end{parts}

\medskip




\noindent   \textbf{Question 20.} In note-taking, what is the Cornell Method?\pts{1}

\begin{parts}
  \item A structured format dividing notes into cues, notes, and summary sections
  \item Writing full paragraphs continuously
  \item Drawing random pictures
  \item Recording audio only
\end{parts}

\medskip




\noindent   \textbf{Question 21.} What is the primary goal of proofreading?\pts{1}

\begin{parts}
  \item Identifying and correcting spelling, punctuation, and grammatical errors
  \item Rewriting the entire text from scratch
  \item Changing font color
  \item Increasing page numbers
\end{parts}

\medskip




\noindent   \textbf{Question 22.} What is a paragraph's ``topic sentence''?\pts{1}

\begin{parts}
  \item The sentence that states the main idea of the paragraph
  \item The last sentence of an essay
  \item A transition phrase between paragraphs
  \item A random supporting detail
\end{parts}

\medskip




\noindent   \textbf{Question 23.} Coherence in writing refers to:\pts{1}

\begin{parts}
  \item Logical flow and smooth connection between sentences and ideas
  \item Using complex vocabulary only
  \item High word count
  \item Formatting margins correctly
\end{parts}

\medskip




\noindent   \textbf{Question 24.} What is plagiarism?\pts{1}

\begin{parts}
  \item Using someone else's work or ideas without proper acknowledgment
  \item Writing an original essay
  \item Citing academic sources correctly
  \item Peer reviewing a classmate's draft
\end{parts}

\medskip




\noindent   \textbf{Question 25.} Which element is essential in a formal report?\pts{1}

\begin{parts}
  \item Title page, executive summary, findings, and recommendations
  \item Personal anecdotes
  \item Fictional story elements
  \item Informal chat transcript
\end{parts}

\medskip




\noindent   \textbf{Question 26.} What is the function of transitional words (   \textit{however, furthermore, consequently}) in writing?\pts{1}

\begin{parts}
  \item To guide the reader and show relationships between ideas
  \item To fill up space
  \item To make sentences longer
  \item To confuse the reader
\end{parts}

\medskip




\noindent   \textbf{Question 27.} In professional communication, what is an inappropriate email signature?\pts{1}

\begin{parts}
  \item Using informal nicknames or inspirational quotes in business emails
  \item Full name and job title
  \item Contact phone number
  \item Organization address
\end{parts}

\medskip




\noindent   \textbf{Question 28.} What is the purpose of an abstract in academic papers?\pts{1}

\begin{parts}
  \item A concise summary of the entire research paper or article
  \item A list of author biography
  \item Detailed raw data tables
  \item Acknowledgments page
\end{parts}

\medskip




\noindent   \textbf{Question 29.} Which voice is generally preferred in objective formal academic writing?\pts{1}

\begin{parts}
  \item Passive voice or third-person objective perspective
  \item First-person subjective perspective
  \item Second-person direct address (``you'')
  \item Conversational dialogue
\end{parts}

\medskip




\noindent   \textbf{Question 30.} What does ``conciseness'' in writing mean?\pts{1}

\begin{parts}
  \item Expressing ideas clearly using as few words as necessary without losing meaning
  \item Writing extremely short sentences only
  \item Skipping important details
  \item Using repetitive words
\end{parts}

\medskip




\noindent   \textbf{Question 31.} Which format is commonly used for formatting business letters?\pts{1}

\begin{parts}
  \item Block format (left-aligned paragraphs without indentation)
  \item Double-spaced poetry format
  \item Right-aligned text format
  \item Center-aligned text format
\end{parts}

\medskip




\noindent   \textbf{Question 32.} In story writing, what is the ``climax''?\pts{1}

\begin{parts}
  \item The turning point or moment of highest emotional tension in the plot
  \item The introduction of characters
  \item The final resolution
  \item The title page
\end{parts}

\medskip




\noindent   \textbf{Question 33.} What is a ``hook'' in the introduction of an essay?\pts{1}

\begin{parts}
  \item An opening sentence designed to capture the reader's attention
  \item The final conclusion statement
  \item A footnote citation
  \item A definition of a common word
\end{parts}

\medskip




\noindent   \textbf{Question 34.} What is the purpose of a bibliography or reference list?\pts{1}

\begin{parts}
  \item To give credit to all sources consulted during research
  \item To increase word count
  \item To list future books to read
  \item To summarize the thesis
\end{parts}

\medskip




\noindent   \textbf{Question 35.} What is ``freewriting''?\pts{1}

\begin{parts}
  \item Writing continuously without worrying about grammar, spelling, or structure to generate ideas
  \item Copying text from a textbook
  \item Writing a polished final draft directly
  \item Writing without a pen
\end{parts}

\medskip




\noindent   \textbf{Question 36.} What is the main characteristic of narrative writing?\pts{1}

\begin{parts}
  \item Telling a story or recounting a sequence of events
  \item Comparing two statistical tables
  \item Describing a scientific apparatus
  \item Formulating an equation
\end{parts}

\medskip




\noindent   \textbf{Question 37.} When writing a letter of complaint, what tone should be maintained?\pts{1}

\begin{parts}
  \item Firm, polite, and objective with specific facts and evidence
  \item Angry and abusive
  \item Sarcastic and threatening
  \item Apologetic and passive
\end{parts}

\medskip




\noindent   \textbf{Question 38.} In e-communication, what does ``netiquette'' refer to?\pts{1}

\begin{parts}
  \item Acceptable rules of polite and professional behavior online
  \item Internet connection speed
  \item Web design techniques
  \item Email encryption software
\end{parts}

\medskip




\noindent   \textbf{Question 39.} What is a mnemonic device in note-taking and learning?\pts{1}

\begin{parts}
  \item A memory aid or technique (such as acronyms) to help remember information
  \item A digital tablet for taking notes
  \item A type of dictionary
  \item A grammatical rule
\end{parts}

\medskip




\noindent   \textbf{Question 40.} What does it mean to ``reply all'' to an email?\pts{1}

\begin{parts}
  \item Sending your response to the sender and every person listed in To and CC fields
  \item Sending an email to your entire contact list
  \item Deleting the email chain
  \item Forwarding to an external party
\end{parts}

\medskip




\noindent   \textbf{Question 41.} Why should you avoid excessive use of jargon/technical words in e-writing for general audiences?\pts{1}

\begin{parts}
  \item It can confuse readers and reduce clarity
  \item It makes writing too easy
  \item It reduces file size
  \item It speeds up internet speed
\end{parts}

\medskip




\noindent   \textbf{Question 42.} What is the purpose of the executive summary in a formal report?\pts{1}

\begin{parts}
  \item To summarize key findings and recommendations for busy executives
  \item To list index terms
  \item To show raw data sheets
  \item To present budget receipts
\end{parts}

\medskip




\noindent   \textbf{Question 43.} Which of the following is the most important characteristic of an effective précis?\pts{1}

\begin{parts}
  \item Brevity, clarity, and faithfulness to the original meaning
  \item Creative additions
  \item Expressing personal bias
  \item Ornate language
\end{parts}

\medskip




\noindent   \textbf{Question 44.} What is mind mapping used for in pre-writing?\pts{1}

\begin{parts}
  \item Visually organizing ideas and brainstorming connections around a central concept
  \item Correcting punctuation errors
  \item Formatting citations
  \item Printing final documents
\end{parts}

\medskip




\noindent   \textbf{Question 45.} When completing a story, which element is most important to consider?\pts{1}

\begin{parts}
  \item Ensuring a logical resolution that ties up main conflict lines
  \item Adding new characters at the last line
  \item Changing the setting abruptly
  \item Leaving sentences incomplete
\end{parts}

\medskip




\noindent   \textbf{Question 46.} In which situation would it be appropriate to use the ``BCC'' option in an email?\pts{1}

\begin{parts}
  \item Sending an email to a large group while protecting recipients' privacy
  \item Emailing your supervisor directly
  \item Submitting an assignment to a teacher
  \item Asking a coworker a question
\end{parts}

\medskip




\noindent   \textbf{Question 47.} Which of the following is an example of a formal writing style?\pts{1}

\begin{parts}
  \item ``The experiment demonstrated a significant increase in efficiency.''
  \item ``Hey guys, the test was totally awesome!''
  \item ``I don't think it's gonna work at all.''
  \item ``LOL that's so cool.''
\end{parts}

\medskip




\noindent   \textbf{Question 48.} What is an email?\pts{1}

\begin{parts}
  \item A method of exchanging digital messages over internet networks
  \item A physical letter sent via postal post
  \item A voice call recording
  \item A printed newspaper article
\end{parts}

\medskip




\noindent   \textbf{Question 49.} Which of the following is an essential step in writing a précis?\pts{1}

\begin{parts}
  \item Reading the original text carefully to grasp the central theme
  \item Changing facts to fit your view
  \item Translating text into another language
  \item Adding fictional details
\end{parts}

\medskip




\noindent   \textbf{Question 50.} When writing a précis, how should you handle the author's examples and anecdotes?\pts{1}

\begin{parts}
  \item Omit unnecessary details, examples, and anecdotes to maintain brevity
  \item Include every example in full detail
  \item Add extra examples of your own
  \item Replace them with personal stories
\end{parts}

\medskip




\noindent   \textbf{Question 51.} Which of the following should you do after writing a précis draft?\pts{1}

\begin{parts}
  \item Compare it with the original text to ensure accuracy and word limit adherence
  \item Throw it away without checking
  \item Double the word count
  \item Add illustrations
\end{parts}

\medskip




\noindent   \textbf{Question 52.} Which of these should be avoided in a good essay?\pts{1}

\begin{parts}
  \item Vague thesis statements, logical fallacies, and informal slang
  \item Clear paragraph transitions
  \item Evidence-backed arguments
  \item Proper introduction and conclusion
\end{parts}

\medskip




\noindent   \textbf{Question 53.} Which of the following is NOT typically part of a standard essay structure?\pts{1}

\begin{parts}
  \item Vague, disconnected topic sentences
  \item Introduction with thesis statement
  \item Body paragraphs with evidence
  \item Conclusion summarizing key points
\end{parts}

\medskip




\noindent   \textbf{Question 54.} What is a ``dangling modifier'' in grammar?\pts{1}

\begin{parts}
  \item A modifier that misleads or lacks a clear word to modify in the sentence
  \item A type of punctuation mark
  \item A strong verb
  \item A transition word
\end{parts}

\medskip




\noindent   \textbf{Question 55.} In business correspondence, what does ``FYI'' stand for?\pts{1}

\begin{parts}
  \item For Your Information
  \item For Your Inspection
  \item From Your Inbox
  \item Final Yearly Income
\end{parts}

\medskip




\noindent   \textbf{Question 56.} What is the purpose of an appendix in a formal report?\pts{1}

\begin{parts}
  \item To include supplementary material (raw data, questionnaires, or detailed tables)
  \item To list main chapters
  \item To serve as the cover page
  \item To state author qualifications
\end{parts}

\medskip




\noindent   \textbf{Question 57.} Which punctuation mark is used to join two independent clauses without a conjunction?\pts{1}

\begin{parts}
  \item Semicolon (;)
  \item Comma (,)
  \item Hyphen (-)
  \item Apostrophe (')
\end{parts}

\medskip




\noindent   \textbf{Question 58.} What is active voice?\pts{1}

\begin{parts}
  \item A sentence structure where the subject performs the action of the verb
  \item A sentence structure where the subject receives the action
  \item A sentence without a verb
  \item A question sentence
\end{parts}

\medskip




\noindent   \textbf{Question 59.} What is the benefit of outlining before writing an essay?\pts{1}

\begin{parts}
  \item Helps structure ideas logically and ensures coherent flow
  \item Eliminates the need for proofreading
  \item Automatically fixes spelling errors
  \item Reduces word count to zero
\end{parts}

\medskip




\noindent   \textbf{Question 60.} What is a prerequisite for effective note-taking during a lecture?\pts{1}

\begin{parts}
  \item Active listening and identifying main ideas/keywords
  \item Writing down every single spoken word verbatim
  \item Sleeping during the presentation
  \item Copying classmate's notes later
\end{parts}

\medskip




\noindent   \textbf{Question 61.} In formal writing, why should contractions (e.g.,   \textit{don't, can't, won't}) be avoided?\pts{1}

\begin{parts}
  \item Contractions represent an informal tone suited for casual communication
  \item They are grammatically illegal
  \item They confuse spellcheckers
  \item They increase word count
\end{parts}

\medskip




\noindent   \textbf{Question 62.} What is the primary purpose of descriptive writing?\pts{1}

\begin{parts}
  \item To create a vivid sensory image in the reader's mind using detailed imagery
  \item To present statistical data
  \item To argue a legal case
  \item To write computer code
\end{parts}

\medskip




\noindent   \textbf{Question 63.} What is the role of a peer reviewer?\pts{1}

\begin{parts}
  \item To provide constructive feedback on clarity, structure, and content of a draft
  \item To assign final grades
  \item To rewrite the paper entirely
  \item To reject the paper without comments
\end{parts}

\medskip




\noindent   \textbf{Question 64.} What is ``brainstorming''?\pts{1}

\begin{parts}
  \item A creative technique for generating a wide list of ideas on a topic freely
  \item A medical test
  \item Formatting a document
  \item Checking dictionary definitions
\end{parts}

\medskip




\noindent   \textbf{Question 65.} What should be avoided while writing a précis?\pts{1}

\begin{parts}
  \item Using complex technical words and adding personal commentary
  \item Retaining original meaning
  \item Summarizing concisely
  \item Presenting key ideas clearly
\end{parts}

\medskip




\noindent   \textbf{Question 66.} Which of the following are characteristics of a formal writing style?\pts{1}

\begin{parts}
  \item Objective tone and avoidance of informal slang/contractions
  \item Excessive use of emojis
  \item Casual conversational slang
  \item Random capitalization
\end{parts}

\medskip




\noindent   \textbf{Question 67.} What is the purpose of peer review in the writing process?\pts{1}

\begin{parts}
  \item To gain fresh perspective and identify areas for improvement in clarity and logic
  \item To copy another student's work
  \item To skip final editing
  \item To change the topic completely
\end{parts}

\medskip




\noindent   \textbf{Question 68.} What is the most appropriate tone for a formal business letter?\pts{1}

\begin{parts}
  \item Professional, respectful, and clear
  \item Overly casual
  \item Aggressive
  \item Humorous and silly
\end{parts}

\medskip




\noindent   \textbf{Question 69.} What is the most important factor in effective note-taking?\pts{1}

\begin{parts}
  \item Capturing essential concepts and key supporting details clearly
  \item Using fancy handwriting
  \item Using expensive notebooks
  \item Writing at maximum speed
\end{parts}

\medskip




\noindent   \textbf{Question 70.} Which of the following is a common mistake in essay writing?\pts{1}

\begin{parts}
  \item Lack of a clear thesis statement and weak paragraph transitions
  \item Writing with clarity
  \item Organizing paragraphs logically
  \item Proofreading before submission
\end{parts}

\vfill
\rule{\linewidth}{0.4pt}

\begin{center}\small
\href{https://bhu-exam.vercel.app/}{Click here to visit our website}\\[4pt]
\href{https://me-aryan.vercel.app/}{Click here to visit our contributor}
\end{center}

\end{document}
